\section{IST-Analyse}

\subsection{Bestehender Avisierungsprozess}

Im betrachteten Ausgangsprozess ordnete die Disposition bereits
in DISPO vorhandene Heizölaufträge einer Tour zu und
legte den vorgesehenen Liefertag sowie ein konkretes
Lieferzeitfenster fest. Die anschließende Abstimmung des
vorgeschlagenen Termins mit dem Kunden erfolgte weitgehend
manuell. Die dafür erforderlichen Auftrags- und
Lieferinformationen mussten für die jeweilige Kontaktaufnahme
zusammengestellt und dem Kunden übermittelt werden.

Die eingehenden Kundenrückmeldungen wurden anschließend manuell
ausgewertet und dem jeweiligen Auftrag zugeordnet. Abhängig von
der Rückmeldung mussten der Bearbeitungsstatus in DISPO
aktualisiert und gegebenenfalls weitere Maßnahmen eingeleitet
werden. Auch noch ausstehende Antworten mussten durch die
Disposition gesondert nachverfolgt werden.

Zwischen der Kundenkommunikation und der Dokumentation im
DISPO bestand keine automatisierte Verbindung. Mit
zunehmender Anzahl parallel bearbeiteter Lieferungen stiegen
daher sowohl der Koordinationsaufwand als auch die Anforderungen
an eine eindeutige Zuordnung und Nachverfolgung der
Kundenreaktionen.

\subsection{Bestehende Kommunikationskanäle}

Die Abstimmung geplanter Liefertermine erfolgte überwiegend
telefonisch. Die Disposition kontaktierte den Kunden und
übermittelte dabei die für die Rückbestätigung relevanten
Informationen zum Auftrag, zum Liefertag und zum vorgesehenen
Lieferzeitfenster.

In einzelnen Fällen wurden auch alternative Kontaktmöglichkeiten
wie E-Mail oder WhatsApp genutzt. Die Auswahl und Verwendung des
jeweiligen Kommunikationswegs erfolgte durch die Disposition und
wurde für jeden Auftrag individuell gehandhabt.

Unabhängig vom verwendeten Kontaktweg wurden die eingehenden
Kundenrückmeldungen manuell ausgewertet, dem entsprechenden
Auftrag zugeordnet und in DISPO dokumentiert. War der
Kunde nicht erreichbar oder blieb eine eindeutige Rückmeldung
aus, wurden weitere Kontaktversuche ebenfalls durch die
Disposition organisiert und nachverfolgt.

\subsection{Schwachstellen und Optimierungspotenziale}

Die wesentlichen Schwachstellen des bestehenden Prozesses lagen
in den Medienbrüchen zwischen Terminplanung,
Kundenkommunikation, Bearbeitung und Dokumentation.
Auftragsinformationen mussten für die Kontaktaufnahme
zusammengestellt und Kundenrückmeldungen anschließend manuell
dem jeweiligen Auftrag zugeordnet werden. Dadurch entstanden
zusätzliche Bearbeitungsschritte sowie ein erhöhtes Risiko für
Verzögerungen, Übertragungsfehler und Missverständnisse.

Statusinformationen wurden zwar in DISPO dokumentiert,
ihre Aktualität hing jedoch von der manuellen Pflege durch die
Disposition ab. Kontaktversuche, Kundenrückmeldungen und noch
offene Abstimmungen waren nicht in einer gemeinsamen
Prozesssicht zusammengeführt. Der aktuelle Stand eines Vorgangs
konnte daher nur anhand der vorhandenen Einträge und der
bisherigen Kommunikation nachvollzogen werden.

Auch nicht erreichbare Kunden und ausbleibende Rückmeldungen
erforderten eine manuelle Kontrolle. Weitere Kontaktversuche
mussten einzeln organisiert und die betroffenen Aufträge
gesondert nachverfolgt werden. Mit zunehmender Anzahl parallel
bearbeiteter Lieferungen stiegen dadurch der zeitliche,
personelle und organisatorische Aufwand.

Optimierungspotenzial bestand vor allem in der Reduzierung
wiederkehrender manueller Arbeitsschritte, einer einheitlicheren
Bearbeitung der Kundenkommunikation und einer verbesserten
Nachvollziehbarkeit des jeweiligen Rückmeldestands. Gleichzeitig
sollte die bestehende fachliche Verantwortung der Disposition
für die Liefer- und Tourenplanung erhalten bleiben.

\subsection{Funktionale Anforderungen}

Aus der Analyse des bestehenden Prozesses sowie den im
Projektverlauf vorgenommenen fachlichen Konkretisierungen und
Erweiterungen wurden die folgenden funktionalen Anforderungen
an die digitale Lieferavisierung abgeleitet:

\begin{table}[htbp]
    \centering
    \caption{Funktionale Anforderungen}
    \label{tab:funktionale-anforderungen}
    \begin{tabularx}{\textwidth}{@{}p{0.10\textwidth}X@{}}
        \toprule
        ID & Beschreibung \\
        \midrule
        A1 &
        Der Avisierungsservice muss die erforderlichen Auftrags-,
        Kunden-, Kontakt- und Lieferdaten über eine geschützte
        DISPO-Schnittstelle übernehmen und dem zugehörigen Unternehmen
        zuordnen. \\

        A2 &
        Für einen Avisierungsvorgang muss eine individuelle
        Bestätigungsanfrage mit einem personalisierten
        Kundenlink erzeugt werden. \\

        A3 &
        Die Bestätigungsanfrage muss abhängig vom hinterlegten
        Kommunikationskanal über E-Mail, SMS oder WhatsApp
        versendet werden können. \\

        A4 &
        Der Kunde muss über den Bestätigungslink die relevanten
        Auftrags- und Lieferinformationen einsehen und den
        vorgeschlagenen Liefertermin bestätigen oder ablehnen
        können. Optional muss eine Kundennachricht übermittelt
        werden können. \\

        A5 &
        Pro Bestätigungsanfrage darf nur eine Kundenreaktion
        statuswirksam verarbeitet werden. Für weitere Reaktionen darf
        ausschließlich die jeweils aktuelle Bestätigungsanfrage gültig
        sein. \\

        A6 &
        Das System muss die für eine Bestätigungsanfrage
        festgelegte Antwortfrist überwachen und eine
        ausbleibende Kundenreaktion als \texttt{NO\_RESPONSE}
        verarbeiten. \\

        A7 &
        Die aus Kundenreaktionen oder Fristabläufen resultierenden
        Status \texttt{CONFIRMED}, \texttt{REJECTED} und
        \texttt{NO\_RESPONSE} müssen persistent gespeichert und über
        die vorgesehene Rückmeldeschnittstelle an DISPO
        übermittelt werden. \\

        A8 &
        Die Disposition muss über ein geschütztes Dashboard
        Touren, Aufträge, Statusanzahlen und offene
        Avisierungsvorgänge einsehen, durchsuchen und filtern
        können. \\

        A9 &
        Die Auftragsdetailansicht muss die aktuelle
        Bestätigungsanfrage, frühere Anfragen und vorhandene
        Kundenreaktionen darstellen. Eine zulässige erneute
        Avisierung muss über einen auswählbaren
        Kommunikationskanal angestoßen werden können. \\

        A10 &
        Mandantenspezifische E-Mail-Einstellungen müssen über
        das Dashboard verwaltet werden können. Zusätzlich müssen
        ein SMTP-Verbindungstest und der Versand einer
        Testnachricht möglich sein. \\

        A11 &
        Für bestätigte Lieferungen muss dem Kunden am Liefertag eine
        Lieferverfolgung bereitgestellt werden, in der die Zielposition
        und die aktuelle, von DISPO bereitgestellte Fahrerposition
        dargestellt werden. \\
        \bottomrule
    \end{tabularx}
\end{table}

\subsection{Nichtfunktionale Anforderungen}

Neben den funktionalen Anforderungen wurden auch nichtfunktionale Anforderungen definiert:

\begin{table}[htbp]
    \centering
    \caption{Nichtfunktionale Anforderungen}
    \label{tab:nichtfunktionale-anforderungen}
    \begin{tabularx}{\textwidth}{@{}p{0.12\textwidth}X@{}}
        \toprule
        ID & Beschreibung \\
        \midrule
        N1 &
        Auftragsdaten, Bestätigungsanfragen und
        Konfigurationen verschiedener Unternehmen müssen
        mandantenbezogen verarbeitet und voneinander getrennt
        werden. \\

        N2 &
        DISPO-, Dashboard- und Kundenschnittstellen müssen
        entsprechend ihrem jeweiligen Zugriffskontext
        geschützt werden. Zugangsdaten und andere sensible
        Konfigurationswerte dürfen nicht ungeschützt
        offengelegt werden. \\

        N3 &
        Zeitabhängige und asynchrone Prozessschritte müssen
        dauerhaft und nachvollziehbar ausgeführt werden.
        Wiederholte Ausführungen dürfen nicht zu widersprüchlichen
        Statusänderungen oder mehrfach verarbeiteten
        Kundenreaktionen führen. \\

        N4 &
        Fehler an externen Kommunikations- und DISPO-Schnittstellen
        müssen kontrolliert behandelt und protokolliert werden.
        Vorübergehend fehlgeschlagene Prozessschritte müssen, soweit
        fachlich zulässig, erneut ausgeführt werden können. \\

        N5 &
        Die Anwendung muss so strukturiert sein, dass externe
        Systeme, Kommunikationskanäle und technische Adapter
        mit möglichst geringen Auswirkungen auf die fachliche
        Anwendungslogik ausgetauscht oder erweitert werden
        können. \\

        N6 &
        Zeitabhängige fachliche Entscheidungen und zentrale
        Systemabläufe müssen reproduzierbar testbar sein. Die
        wesentlichen fachlichen, integrativen und
        systemübergreifenden Abläufe müssen durch automatisierte Tests
        überprüft werden können. \\

        N7 &
        Personenbezogene und sicherheitsrelevante Daten müssen
        bei Speicherung, Übertragung und Protokollierung
        angemessen geschützt werden. Dynamische Tracking-Daten
        dürfen nicht als dauerhaft gültige Antworten
        zwischengespeichert werden. \\

        N8 &
        Die kunden- und dispositionsseitigen Oberflächen müssen
        die jeweils relevanten Informationen verständlich
        darstellen und eine eindeutige Bedienung der
        vorgesehenen Aktionen ermöglichen. \\
        \bottomrule
    \end{tabularx}
\end{table}